% ============================================================
% The Interaction Gap
% ============================================================

\documentclass{article}

% --- Packages ---
\usepackage[utf8]{inputenc}
\usepackage[T1]{fontenc}
\usepackage{amsmath,amssymb,amsfonts}
\usepackage{graphicx}
\usepackage{booktabs}
\usepackage{hyperref}
\usepackage{xcolor}
\usepackage{multirow}
\usepackage{enumitem}
\usepackage{natbib}
\usepackage[margin=1in]{geometry}
\usepackage{thmtools}

% --- Theorem environments ---
\newtheorem{definition}{Definition}

% --- Hyperlink styling ---
\hypersetup{
    colorlinks=true,
    linkcolor=blue!70!black,
    citecolor=blue!70!black,
    urlcolor=blue!70!black
}

% --- Title ---
\title{%
  \textbf{The Interaction Gap: Why Text Embedding Benchmarks\\Miss What Matters}
}

\author{
  Anonymous Authors
}

\date{}

% ============================================================
\begin{document}
\maketitle

% ============================================================
% ABSTRACT
% ============================================================
\begin{abstract}
The Massive Text Embedding Benchmark (MTEB) is the de facto standard for evaluating text embedding models. All eight of its task types evaluate \emph{similarity}---whether two texts are topically related, semantically close, or mutually relevant. None evaluate \emph{interaction}---whether two texts produce a specific functional outcome when combined, such as entailment, contradiction, or compatibility.

We demonstrate this gap using MTEB's own data. The SICK-R dataset is evaluated by MTEB for semantic similarity (relatedness scores 1--5) and by us for textual entailment (entailment/contradiction/neutral labels)---both annotations exist on every sentence pair. Across 15 embedding models spanning the full MTEB quality range (Spearman $\rho$ 0.55--0.81 on STS), we find that MTEB rankings are \textbf{inversely correlated} with entailment separation performance: $\rho = -0.761$, $p < 0.001$. The model ranked highest by MTEB is among the worst at distinguishing entailment from contradiction. The only model that correctly separates them is GloVe---the one MTEB ranks last.

This inversion is structural, not accidental. Contradictory sentence pairs in SICK-R have \emph{higher} human relatedness ratings (mean 4.57/5) than entailing pairs (3.88/5), because a statement and its negation are about the same topic. A model optimized for relatedness will therefore assign high similarity to contradictions---exactly backwards for entailment. We complement this finding with experiments on FEVER claim verification (9 models) and SummEval summary quality (7 models), showing that MTEB predictions are accurate when similarity and interaction align but fail when they conflict. These results reveal a structural blind spot in embedding evaluation with implications for fact-checking, entailment-based retrieval, and compatibility scoring.
\end{abstract}

% ============================================================
% 1. INTRODUCTION
% ============================================================
\section{Introduction}
\label{sec:introduction}

Text embeddings are the backbone of modern information retrieval, semantic search, and retrieval-augmented generation \citep{karpukhin2020dpr, lewis2020rag}. When practitioners need to select an embedding model, they overwhelmingly turn to the Massive Text Embedding Benchmark (MTEB) leaderboard \citep{muennighoff2023mteb}, which ranks models across 56 datasets spanning 8 task types. The implicit promise is comprehensive: a model that ranks well overall should perform well on your task.

We show this promise has a structural blind spot. All eight MTEB task types---Bitext Mining, Classification, Clustering, Pair Classification, Reranking, Retrieval, Semantic Textual Similarity (STS), and Summarization---evaluate whether embeddings capture \emph{similarity} between texts. They ask: are these texts about the same thing? Is this document relevant to this query? How semantically close are these sentences?

But many real-world applications require a different judgment:

\begin{enumerate}[leftmargin=*, itemsep=2pt]
    \item \textbf{Claim verification.} A fact-checking system retrieves documents relevant to a claim. Both supporting and refuting documents are ``relevant'' (high similarity), but the system must distinguish their \emph{interaction valence}---support versus refutation.

    \item \textbf{Textual entailment.} A legal system must determine whether clause A entails clause B. The clauses may share vocabulary and topic (high similarity) yet have opposite entailment relations.

    \item \textbf{Compatibility scoring.} A matching system must predict whether two entities are \emph{compatible}---a relationship that is neither symmetric nor monotonic in similarity. Two job candidates with identical resumes are maximally similar but not both optimal for the same role.
\end{enumerate}

We call the gap between what MTEB measures and what these applications require the \textbf{interaction gap}.

The most striking aspect of this gap is that it hides in plain sight. The SICK-R dataset \citep{marelli2014sick}---already used by MTEB for semantic similarity evaluation---contains entailment labels on every sentence pair. MTEB evaluates the similarity labels and ignores the entailment labels. We evaluate both and show the rankings invert.

Our contributions:

\begin{enumerate}[leftmargin=*, itemsep=2pt]
    \item We formally distinguish \emph{similarity tasks} from \emph{interaction tasks} and show why cosine similarity is structurally limited for the latter (\S\ref{sec:problem}).

    \item We evaluate 15 embedding models on the same SICK-R data used by MTEB but with an interaction question, finding a highly significant inverse correlation between MTEB and interaction rankings: $\rho = -0.761$, $p < 0.001$ (\S\ref{sec:results}).

    \item We show this inversion is structural---it arises because contradictions are topically related---and demonstrate on two additional tasks (FEVER, SummEval) that the gap appears precisely when similarity and interaction conflict (\S\ref{sec:analysis}).

    \item We release our evaluation code as an open-source benchmark to enable interaction-aware embedding evaluation alongside MTEB.
\end{enumerate}

% ============================================================
% 2. RELATED WORK
% ============================================================
\section{Related Work}
\label{sec:related}

\textbf{Embedding benchmarks.} MTEB \citep{muennighoff2023mteb} standardized embedding evaluation across 8 task types and 56 datasets, building on BEIR's zero-shot retrieval evaluation \citep{thakur2021beir}. MTEB v2 (October 2025) added multimodal and cross-encoder support but no interaction task types. MMTEB \citep{mmteb2025} expanded to 500+ tasks across 250+ languages, again without interaction tasks. The MTEB maintainers themselves acknowledge gaps in domain coverage and the limitations of cosine similarity as a universal metric \citep{chung2025maintaining}. PTEB \citep{pteb2025} demonstrated that MTEB scores are inflatable via overfitting to static test sets, further questioning the benchmark's reliability as a comprehensive quality signal.

\textbf{Cosine similarity limitations.} \citet{steck2024cosine} proved analytically that cosine similarity on learned embeddings can yield arbitrary results due to gauge freedom in training. \citet{you2025semantics} cataloged four failure modes including the inability to distinguish entailment strength: ``entailment and NLI tasks may produce indistinguishable similarity scores for weak and strong hypotheses.'' \citet{ethayarajh2019contextual} showed that contextual embeddings suffer from anisotropy, concentrating vectors in narrow cones that inflate similarity scores.

\textbf{Theoretical embedding limitations.} \citet{weller2025limit} proved that the number of distinct top-$k$ retrieval subsets encodable by $d$-dimensional embeddings is fundamentally bounded. Their LIMIT dataset---50K documents describing people and preferences, with queries asking ``who likes X?''---showed embedding models achieving only 12.9\% recall while sparse retrieval (BM25) achieved 93.6\%. Performance on LIMIT was uncorrelated with BEIR/MTEB scores, providing theoretical grounding for why similarity-trained embeddings fail on interaction tasks.

\textbf{STS does not predict downstream performance.} \citet{abe2022sts} directly demonstrated that STS benchmark performance does not reliably predict downstream task performance, confirming that similarity evaluation is an incomplete proxy for general embedding quality.

\textbf{Asymmetric representations.} The inadequacy of symmetric similarity for asymmetric relations is well-studied. \citet{vilnis2015word} proposed Gaussian embeddings where KL divergence naturally captures asymmetric entailment. \citet{vendrov2016order} introduced order embeddings with partial-order structure for entailment. \citet{nickel2017poincare} embedded hierarchies in hyperbolic space. \citet{yoda2024gausscse} extended Gaussian representations to sentences, representing entailment as containment between regions. These approaches share the insight that point-vector cosine similarity cannot capture directional or asymmetric relations.

\textbf{Similarity vs.\ complementarity.} \citet{mcauley2015complementary} showed that product substitutability and complementarity require different embedding spaces---a single cosine similarity cannot capture both. \citet{kvernadze2022dual} used dual embeddings where the IN-IN dot product captures similarity and the IN-OUT dot product captures complementarity. The distinction between similarity and interaction parallels the distinction between substitutes and complements: compatible entities are often dissimilar, and similar entities are often not complementary.

\textbf{Cross-encoders.} Bi-encoders encode each text independently and compute similarity via dot product; cross-encoders jointly encode both texts and compute an interaction score \citep{nogueira2019passage, humeau2020polyencoders}. Cross-encoders consistently outperform bi-encoders on tasks requiring fine-grained inter-input reasoning \citep{thakur2021beir}. Our work does not compare bi-encoders to cross-encoders---that gap is well-established. Instead, we show that \emph{the evaluation of bi-encoders themselves} has a blind spot: MTEB measures only the similarity dimension of bi-encoder quality, missing the interaction dimension entirely.

\textbf{NLI for embedding training.} Sentence-BERT \citep{reimers2019sentence} and SimCSE \citep{gao2021simcse} use NLI data to train embeddings, showing that entailment signal improves similarity performance. Critically, their NLI-trained models are still \emph{evaluated} on similarity tasks (STS). Our work asks the converse: when evaluated on entailment, do similarity-optimized models perform well? Our results show they do not---in fact, they perform worse than untrained baselines.

% ============================================================
% 3. THE PROBLEM: SIMILARITY vs. INTERACTION
% ============================================================
\section{Similarity vs.\ Interaction}
\label{sec:problem}

\begin{definition}[Similarity Task]
\label{def:similarity}
A \emph{similarity task} over text pairs $(a, b)$ is one where the target label $y$ is monotonically increasing in the semantic relatedness between $a$ and $b$. More related texts always receive higher (or equal) scores.
\end{definition}

\begin{definition}[Interaction Task]
\label{def:interaction}
An \emph{interaction task} over text pairs $(a, b)$ is one where the target label $y$ depends on a functional relationship between $a$ and $b$ that is \textbf{not} monotonic in their semantic relatedness. Two pairs with identical relatedness may have different---even opposite---interaction labels.
\end{definition}

In a similarity task, knowing that two texts are ``more related'' always provides information about the label. In an interaction task, relatedness is necessary but not sufficient: you must also understand \emph{how} the texts interact.

\subsection{Why Cosine Similarity Is Structurally Limited}

A bi-encoder maps texts to $d$-dimensional embeddings via $f: \mathcal{T} \to \mathbb{R}^d$ and scores pairs by:
\begin{equation}
    s(a, b) = \cos(f(a), f(b)) = \frac{f(a) \cdot f(b)}{\|f(a)\| \|f(b)\|}
\end{equation}

This has two properties that conflict with interaction tasks:

\textbf{Symmetry.} $s(a, b) = s(b, a)$, but entailment is directional: ``all cats are animals'' does not imply ``all animals are cats.''

\textbf{Monotonicity in relatedness.} Cosine similarity increases as embeddings become more aligned. But interaction tasks require \emph{discrimination within the high-similarity region}. A claim and its supporting evidence have high cosine similarity. A claim and its refuting evidence \emph{also} have high cosine similarity---they share topic, entities, and vocabulary. The interaction label (support vs.\ refute) varies while similarity remains roughly constant.

\subsection{The SICK-R Paradox}

The SICK-R dataset provides the cleanest illustration. Each sentence pair has both a relatedness score (1--5) and an entailment label (E/C/N). The mean relatedness by label is:

\begin{center}
\begin{tabular}{lcc}
\toprule
\textbf{Label} & \textbf{Mean Relatedness} & \textbf{Count} \\
\midrule
Contradiction & 4.57 & 1,424 \\
Entailment & 3.88 & 2,821 \\
Neutral & 2.99 & 5,682 \\
\bottomrule
\end{tabular}
\end{center}

Contradictions have the \emph{highest} relatedness. This is linguistically natural: ``A man is playing guitar'' and ``No man is playing guitar'' are about the same scene---they differ only in polarity. They are maximally related (about the same thing) and maximally opposed (opposite truth values).

A model optimized to predict relatedness scores (MTEB's STS objective) will learn: contradictions $>$ entailments $>$ neutral in similarity. But the correct interaction ranking is: entailments $>$ neutral $>$ contradictions. \textbf{The optimization objectives are in direct conflict.}

This is not a quirk of SICK-R. Any domain where semantic opposition co-occurs with topical overlap will exhibit the same paradox. Fact-checking (claims and refutations share topic), legal analysis (opposing interpretations share terminology), and compatibility scoring (incompatible candidates share demographics) all have this structure.

% ============================================================
% 4. EXPERIMENTAL SETUP
% ============================================================
\section{Experimental Setup}
\label{sec:experiments}

\subsection{Models}

We evaluate 15 embedding models spanning the full MTEB quality range, from static word vectors to state-of-the-art transformer encoders. Models are selected to cover weak, medium, and strong MTEB performance:

\begin{table}[h]
\centering
\caption{Models evaluated, ordered by MTEB STS Spearman on SICK-R. The range spans 0.554--0.812, covering models from static word vectors to the top of the MTEB STS leaderboard.}
\label{tab:models}
\small
\begin{tabular}{@{}llcc@{}}
\toprule
\textbf{Model} & \textbf{Type} & \textbf{Params} & \textbf{MTEB STS $\rho$} \\
\midrule
GloVe 6B 300d & Static word vectors & -- & 0.554 \\
paraphrase-MiniLM-L3-v2 & Transformer & 17M & 0.762 \\
all-MiniLM-L6-v2 & Transformer & 22M & 0.772 \\
e5-base-v2 & Transformer & 109M & 0.775 \\
nomic-embed-text-v1.5 & Transformer & 137M & 0.780 \\
gte-base & Transformer & 109M & 0.783 \\
e5-large-v2 & Transformer & 335M & 0.787 \\
e5-small-v2 & Transformer & 33M & 0.787 \\
gte-large & Transformer & 335M & 0.791 \\
bge-small-en-v1.5 & Transformer & 33M & 0.791 \\
multilingual-e5-large & Transformer & 560M & 0.798 \\
bge-base-en-v1.5 & Transformer & 109M & 0.799 \\
all-distilroberta-v1 & Transformer & 82M & 0.801 \\
all-mpnet-base-v2 & Transformer & 109M & 0.805 \\
bge-large-en-v1.5 & Transformer & 335M & 0.812 \\
\bottomrule
\end{tabular}
\end{table}

\subsection{Tasks}

We evaluate three tasks spanning different interaction types. All use existing, publicly available datasets.

\subsubsection{SICK-R Entailment (Primary)}

The SICK-R dataset \citep{marelli2014sick} contains 4,906 sentence pairs, each annotated with both a relatedness score (1--5) and an entailment label (Entailment, Contradiction, Neutral). MTEB evaluates Spearman correlation between cosine similarity and relatedness scores. We evaluate whether cosine similarity separates entailment from contradiction.

\textbf{Metric: Separation score} (Cohen's $d$). For each model, we compute the mean cosine similarity for entailment pairs ($\bar{s}_E$) and contradiction pairs ($\bar{s}_C$), then:
\begin{equation}
    d = \frac{\bar{s}_E - \bar{s}_C}{s_{\text{pooled}}}
\end{equation}
where $s_{\text{pooled}}$ is the pooled standard deviation. A positive $d$ indicates correct separation (entailment $>$ contradiction). A negative $d$ indicates inverted separation.

\subsubsection{FEVER Claim Verification}

The FEVER dataset \citep{thorne2018fever} contains claims paired with evidence passages labeled SUPPORTS or REFUTES. MTEB uses FEVER for retrieval evaluation (nDCG@10). We evaluate whether cosine similarity between claim and evidence discriminates support from refutation.

\textbf{Metrics:} Separation score (Cohen's $d$: $\bar{s}_{\text{support}} - \bar{s}_{\text{refute}}$) and AUROC for binary classification.

\subsubsection{SummEval Summary Quality}

SummEval \citep{fabbri2021summeval} contains 1,600 source-summary pairs (100 sources $\times$ 16 summaries) with human quality ratings. We evaluate whether cosine similarity between source and summary predicts summary quality.

\textbf{Metric:} Mean Spearman correlation between cosine similarity ranking and human quality ranking, computed per source and averaged.

\subsection{Rank Correlation Analysis}

For each task, we rank all models by their MTEB STS score and by their interaction task score, then compute Spearman's rank correlation $\rho$ between the two rankings. Low or negative $\rho$ indicates that MTEB rankings do not predict interaction performance. We focus on the \textbf{same-dataset comparison} for SICK-R, where both the MTEB and interaction evaluations use identical sentence pairs.

% ============================================================
% 5. RESULTS
% ============================================================
\section{Results}
\label{sec:results}

\subsection{SICK-R: Same Data, Inverted Rankings}

Table~\ref{tab:sickr_results} presents the headline result. Across 15 models, the Spearman rank correlation between MTEB STS performance and entailment separation is $\rho = -0.761$ ($p < 0.001$). The rankings are not merely uncorrelated---they are \textbf{inversely correlated}. Improving at MTEB's similarity task actively \emph{harms} entailment discrimination.

\begin{table}[t]
\centering
\caption{SICK-R results: same 4,906 sentence pairs, same models, different evaluation question. MTEB evaluates Spearman correlation with relatedness scores. We evaluate whether cosine similarity separates entailment from contradiction (Cohen's $d$; positive = correct). Models sorted by MTEB score. \textbf{Rank correlation: $\rho = -0.761$, $p < 0.001$.}}
\label{tab:sickr_results}
\small
\begin{tabular}{@{}lcccccc@{}}
\toprule
\textbf{Model} & \textbf{MTEB} & \textbf{MTEB} & \textbf{Separation} & \textbf{MITE} & $\bar{s}_E$ & $\bar{s}_C$ \\
 & \textbf{STS $\rho$} & \textbf{Rank} & \textbf{(Cohen's $d$)} & \textbf{Rank} & & \\
\midrule
GloVe 6B 300d       & 0.554 & 15 & \textbf{+0.530} & \textbf{1} & 0.951 & 0.897 \\
paraphrase-MiniLM-L3 & 0.762 & 14 & $-0.959$ & 4 & 0.765 & 0.876 \\
all-MiniLM-L6-v2     & 0.772 & 13 & $-0.902$ & 3 & 0.757 & 0.864 \\
e5-base-v2           & 0.775 & 12 & $-1.418$ & 6 & 0.893 & 0.950 \\
nomic-embed-v1.5     & 0.780 & 11 & $-0.999$ & 5 & 0.852 & 0.922 \\
gte-base             & 0.783 & 10 & $-1.078$ & 5 & 0.923 & 0.959 \\
e5-large-v2          & 0.787 & 8 & $-2.256$ & 13 & 0.870 & 0.956 \\
e5-small-v2          & 0.787 & 9 & $-2.169$ & 12 & 0.882 & 0.960 \\
gte-large            & 0.791 & 7 & $-1.280$ & 7 & 0.919 & 0.961 \\
bge-small-en-v1.5    & 0.791 & 6 & $-1.882$ & 10 & 0.776 & 0.918 \\
multilingual-e5-large & 0.798 & 5 & $-2.299$ & 14 & 0.873 & 0.958 \\
bge-base-en-v1.5     & 0.799 & 4 & $-1.860$ & 9 & 0.762 & 0.909 \\
all-distilroberta-v1 & 0.801 & 3 & $-1.856$ & 8 & 0.629 & 0.858 \\
all-mpnet-base-v2    & 0.805 & 2 & $-2.456$ & \textbf{15} & 0.545 & 0.860 \\
bge-large-en-v1.5    & 0.812 & 1 & $-2.108$ & 11 & 0.751 & 0.914 \\
\bottomrule
\end{tabular}
\end{table}

Three observations stand out:

\textbf{GloVe is the only model that gets it right.} With a separation of $+0.530$, GloVe~\citep{pennington2014glove} is the only model where entailment pairs have higher cosine similarity than contradiction pairs. GloVe was not trained on similarity objectives---it captures word co-occurrence statistics that happen to preserve some entailment signal. Every transformer-based model, trained explicitly on similarity, gets the separation backwards.

\textbf{The best MTEB model is nearly the worst.} all-mpnet-base-v2, ranked \#2 on MTEB STS (0.805), has the worst entailment separation ($-2.456$). bge-large-en-v1.5, ranked \#1 (0.812), has the third-worst ($-2.108$). The models MTEB recommends most strongly are the ones most confidently wrong about entailment.

\textbf{Every transformer model assigns higher similarity to contradictions than entailments.} All 14 transformer models have $\bar{s}_C > \bar{s}_E$, confirming the structural prediction from \S\ref{sec:problem}: models trained to predict relatedness learn that contradictions (relatedness 4.57) should score higher than entailments (3.88). The better they learn this, the more inverted their entailment judgments become.

\subsection{FEVER: When Similarity and Interaction Align}

Table~\ref{tab:fever_results} shows FEVER results for 9 models. The rank correlation between MTEB STS and FEVER interaction scores is $\rho = +0.617$ ($p = 0.077$)---positive and approaching significance.

\begin{table}[t]
\centering
\caption{FEVER claim verification results (9 models, 968 balanced claim-evidence pairs). Unlike SICK-R, FEVER shows \emph{positive} (though non-significant) correlation between MTEB and interaction scores. All models successfully separate SUPPORTS from REFUTES. \textbf{Rank correlation: $\rho = +0.617$, $p = 0.077$.}}
\label{tab:fever_results}
\small
\begin{tabular}{@{}lcccc@{}}
\toprule
\textbf{Model} & \textbf{MTEB STS $\rho$} & \textbf{Separation} & \textbf{AUROC} & $\bar{s}_S - \bar{s}_R$ \\
\midrule
bge-base-en-v1.5    & 0.799 & 0.949 & 0.750 & +0.075 \\
bge-small-en-v1.5   & 0.791 & 0.917 & 0.748 & +0.069 \\
e5-small-v2         & 0.787 & 0.852 & 0.727 & +0.031 \\
all-mpnet-base-v2   & 0.805 & 0.815 & 0.723 & +0.105 \\
all-MiniLM-L6-v2    & 0.772 & 0.786 & 0.712 & +0.096 \\
e5-base-v2          & 0.775 & 0.758 & 0.708 & +0.028 \\
e5-large-v2         & 0.787 & 0.731 & 0.697 & +0.028 \\
nomic-embed-v1.5    & 0.780 & 0.689 & 0.689 & +0.048 \\
gte-base            & 0.783 & 0.680 & 0.688 & +0.020 \\
\bottomrule
\end{tabular}
\end{table}

This result is informative \emph{because} it contrasts with SICK-R. In FEVER, supporting evidence is naturally more similar to a claim than refuting evidence---they share factual content, not just topic. There is no structural paradox: similarity and interaction are \emph{aligned}. The MTEB-trained models' ability to detect similarity partially transfers to detecting the support/refute distinction.

\subsection{SummEval: Weak Positive Signal}

For 7 models evaluated on SummEval, the rank correlation between MTEB STS and summary quality discrimination is $\rho = +0.500$ ($p = 0.253$)---positive but not significant with limited models. Summary quality is partially a similarity task: summaries that resemble their source tend to score higher on relevance (though not necessarily on coherence or consistency). Consistent with our framework, MTEB partially predicts performance when similarity and interaction are partially aligned.

\subsection{The Pattern: Alignment Determines Predictability}

Table~\ref{tab:pattern} summarizes across tasks.

\begin{table}[t]
\centering
\caption{MTEB's predictive value depends on whether similarity and interaction are aligned or opposed. When they conflict (SICK-R), MTEB rankings are inversely correlated with interaction performance. When they align (FEVER, SummEval), MTEB partially predicts interaction.}
\label{tab:pattern}
\small
\begin{tabular}{@{}lllcc@{}}
\toprule
\textbf{Task} & \textbf{Alignment} & \textbf{Why} & $\rho$ & $p$ \\
\midrule
SICK-R Entailment & \textbf{Opposed} & Contradictions are more related & $-0.761$ & $<0.001$ \\
FEVER Verification & Aligned & Supporting evidence is more similar & $+0.617$ & $0.077$ \\
SummEval Quality & Partially aligned & Good summaries resemble source & $+0.500$ & $0.253$ \\
\bottomrule
\end{tabular}
\end{table}

The conclusion is not that MTEB is wrong, but that \textbf{MTEB is incomplete}. It accurately evaluates the similarity dimension of embedding quality but is blind to the interaction dimension. When the downstream task aligns similarity with interaction, MTEB rankings transfer. When they conflict---as they structurally must in entailment, fact-checking polarity, and compatibility scoring---MTEB rankings actively mislead.

% ============================================================
% 6. ANALYSIS
% ============================================================
\section{Analysis}
\label{sec:analysis}

\subsection{Why the Inversion Is Structural}

The SICK-R inversion arises from a necessary property of natural language: semantic opposition requires semantic overlap. You can only contradict a claim about the same topic. ``The cat is sleeping'' and ``The cat is not sleeping'' share every content word---they differ only in polarity. This makes their semantic relatedness high (4.57/5 in SICK-R) while their interaction (entailment vs.\ contradiction) is maximally opposed.

Any model trained to optimize relatedness prediction will learn this pattern and assign high similarity to contradictions. The better it learns MTEB's objective, the more it confuses contradiction with entailment. This is not a failure of training---it is the \emph{correct} outcome of the similarity objective. The model is doing exactly what MTEB asks: predicting relatedness. The problem is that MTEB asks the wrong question for interaction tasks.

\subsection{Why GloVe Succeeds}

GloVe \citep{pennington2014glove} is the only model with positive separation. This is counterintuitive---GloVe is a static word-vector average with no sentence-level training. Why does it outperform all transformer models?

GloVe embeddings are trained on word co-occurrence, not on sentence-pair similarity judgments. They have not learned the SICK-R relatedness distribution and therefore have not learned that contradictions should score higher than entailments. Their co-occurrence statistics happen to preserve some entailment signal: entailing sentences share more context words (``A man is playing guitar'' $\to$ ``Someone is making music'') than contradicting sentences (``A man is playing guitar'' $\to$ ``No man is playing guitar''), even though contradicting sentences share more \emph{surface} words.

This finding parallels the ``unlearning'' phenomenon in machine learning: sometimes training on the wrong objective produces worse results than no training at all. Transformer models are not simply ignorant of entailment---they have been trained to systematically mispredict it.

\subsection{Interaction as a Spectrum}

Not all interaction tasks are equally affected by the similarity-interaction gap. Our three tasks fall on a spectrum:

\textbf{Maximally conflicting:} SICK-R entailment. Contradictions have \emph{higher} relatedness than entailments. Similarity optimization directly opposes interaction accuracy. MTEB is anti-predictive ($\rho = -0.761$).

\textbf{Aligned:} FEVER claim verification. Supporting evidence is more similar to claims than refuting evidence. Similarity optimization partially helps interaction accuracy. MTEB is partially predictive ($\rho = +0.617$).

\textbf{Partially aligned:} SummEval quality. Good summaries tend to resemble their source, but quality also depends on coherence and consistency, which are not similarity. MTEB is weakly predictive ($\rho = +0.500$).

This spectrum has practical implications. Practitioners should not assume their interaction task falls at the ``aligned'' end. Any task where semantic opposition co-occurs with topical overlap---entailment, fact-checking polarity, argumentation, legal contradiction, compatibility scoring---is at risk for the similarity-interaction inversion.

\subsection{Connection to Theoretical Results}

Our empirical findings are consistent with the theoretical results of \citet{weller2025limit}, who proved fundamental bounds on the expressiveness of fixed-dimensional embeddings for combinatorial retrieval tasks. Their LIMIT dataset---where queries ask ``who likes X?'' about person descriptions---requires the kind of interaction reasoning that our SICK-R entailment task does: determining a functional relationship (preference match) between two texts, not their topical similarity. The catastrophic failure of all embedding models on LIMIT (12.9\% recall vs.\ 93.6\% for BM25) is the extreme case of the interaction gap we measure.

\citet{steck2024cosine}'s proof that cosine similarity can be arbitrary due to gauge freedom provides another theoretical basis: even if a model \emph{could} encode interaction information, cosine similarity might not surface it. Our results show this is not merely theoretical---in the concrete case of SICK-R, cosine similarity systematically surfaces the \emph{wrong} information for the interaction task.

% ============================================================
% 7. DISCUSSION
% ============================================================
\section{Discussion}
\label{sec:discussion}

\subsection{Implications for Practitioners}

\textbf{Model selection.} Practitioners building systems that require interaction judgments---entailment detection, claim polarity classification, compatibility scoring---should not rely on MTEB rankings alone. Our results show that MTEB's top-ranked models can be the worst choices for these tasks. Task-specific evaluation on interaction benchmarks is necessary.

\textbf{System design.} A retrieval system that uses cosine similarity to find evidence for fact-checking will retrieve both supporting and refuting evidence with similar scores, regardless of which embedding model is used. The choice of embedding model (and its MTEB rank) is irrelevant to this failure mode. Practitioners need post-retrieval interaction classifiers, cross-encoders, or embedding models specifically trained for interaction discrimination---none of which are evaluated by MTEB.

\textbf{Evaluation practice.} We recommend that embedding evaluation suites include at least one interaction task where similarity and interaction are known to conflict. SICK-R entailment serves this purpose using data MTEB already includes, requiring no additional data collection.

\subsection{Relationship to nDCG@10}

Some practitioners have moved toward nDCG@10 as a primary evaluation metric, following the practice of major embedding providers. nDCG@10 is a retrieval ranking metric that is more practically relevant than STS Spearman for search and RAG applications. However, nDCG@10 is still a similarity/relevance metric---it measures whether relevant documents are ranked highly, not whether the model can distinguish types of relevance (e.g., supporting vs.\ refuting evidence). The interaction gap is orthogonal to the choice between STS and nDCG@10: both are blind to interaction valence.

\subsection{Limitations}

\textbf{MTEB score proxy.} We use each model's STS Spearman on SICK-R as the MTEB comparison score, not the model's overall MTEB leaderboard rank. This is the most controlled comparison (same data, different question) but does not directly evaluate whether overall MTEB rank is predictive. Future work should include official leaderboard scores.

\textbf{Model coverage.} Our 15 models span MTEB STS scores from 0.554 to 0.812 but do not include recent frontier models (e.g., those from OpenAI, Cohere, Voyage) or models with MTEB scores below 0.50 or above 0.85. Extending the model set could strengthen or weaken the observed correlations.

\textbf{Three tasks.} We evaluate three interaction tasks. Important interaction types---including pragmatic implication, argumentation strength, and interpersonal compatibility---are not yet covered. SICK-R provides the cleanest same-data comparison; additional tasks would strengthen the generality claim.

\textbf{English only.} All experiments are in English. The interaction gap may differ in other languages, particularly those with different negation or contradiction patterns.

\textbf{Cosine only.} We restrict to cosine similarity to match MTEB's evaluation protocol. Alternative similarity functions (e.g., learned metrics, asymmetric functions) might recover interaction signal from the same embeddings.

% ============================================================
% 8. CONCLUSION
% ============================================================
\section{Conclusion}
\label{sec:conclusion}

We have identified a structural blind spot in text embedding evaluation. The Massive Text Embedding Benchmark (MTEB), despite evaluating 56 datasets across 8 task types, measures only one axis of embedding quality: similarity. Interaction tasks---where the functional relationship between texts matters beyond topical relatedness---are absent.

Using MTEB's own SICK-R dataset, we show that MTEB rankings are \emph{inversely correlated} with entailment discrimination across 15 models ($\rho = -0.761$, $p < 0.001$). This inversion is structural: contradictions are topically related, so models optimized for relatedness systematically confuse contradiction with entailment. The only model that correctly separates them is GloVe---an untrained word-vector average that MTEB ranks last.

The interaction gap is not a universal failure of MTEB. On tasks where similarity and interaction align (FEVER claim verification, summary quality), MTEB rankings partially transfer. The gap matters precisely when similarity and interaction conflict---a condition that arises structurally in entailment, fact-checking polarity, legal contradiction, and compatibility scoring.

We propose that embedding evaluation suites include interaction tasks alongside similarity tasks. The data already exists in MTEB's datasets---it only needs to be evaluated. Our code is publicly available to enable this evaluation.

The interaction gap is not a flaw in any particular model. It is a gap in how we evaluate models. As embedding models are deployed in increasingly complex reasoning systems, closing this gap becomes essential.

% ============================================================
% REFERENCES
% ============================================================
\bibliographystyle{plainnat}
\bibliography{references}

% ============================================================
% APPENDIX
% ============================================================
\appendix

\section{SICK-R Relatedness Distribution}
\label{app:sickr}

The SICK-R dataset contains 9,927 sentence pairs with dual annotations: relatedness scores (1--5) and entailment labels (Entailment/Contradiction/Neutral). We use the 4,906-pair test split for evaluation.

The distribution of relatedness scores by entailment label reveals the structural conflict between similarity and interaction evaluation:

\begin{center}
\begin{tabular}{@{}lccc@{}}
\toprule
\textbf{Entailment Label} & \textbf{Mean Relatedness} & \textbf{Std.} & \textbf{Count} \\
\midrule
Contradiction & 4.57 & 0.49 & 1,424 \\
Entailment & 3.88 & 0.72 & 2,821 \\
Neutral & 2.99 & 0.91 & 5,682 \\
\bottomrule
\end{tabular}
\end{center}

The ordering---Contradiction $>$ Entailment $>$ Neutral in relatedness---is the inverse of what an entailment model should produce. A model that perfectly predicts relatedness scores will assign highest similarity to contradictions, systematically misclassifying them.

\section{Full Model Results}
\label{app:full_results}

\subsection{SICK-R: All 15 Models}

Complete results including mean cosine similarity for each entailment class:

\begin{center}
\small
\begin{tabular}{@{}lccccc@{}}
\toprule
\textbf{Model} & \textbf{MTEB $\rho$} & \textbf{Sep. ($d$)} & $\bar{s}_E$ & $\bar{s}_C$ & $\bar{s}_N$ \\
\midrule
GloVe 6B 300d & 0.554 & +0.530 & 0.951 & 0.897 & -- \\
paraphrase-MiniLM-L3 & 0.762 & $-0.959$ & 0.765 & 0.876 & -- \\
all-MiniLM-L6-v2 & 0.772 & $-0.902$ & 0.757 & 0.864 & -- \\
e5-base-v2 & 0.775 & $-1.418$ & 0.893 & 0.950 & -- \\
nomic-embed-v1.5 & 0.780 & $-0.999$ & 0.852 & 0.922 & -- \\
gte-base & 0.783 & $-1.078$ & 0.923 & 0.959 & -- \\
e5-large-v2 & 0.787 & $-2.256$ & 0.870 & 0.956 & -- \\
e5-small-v2 & 0.787 & $-2.169$ & 0.882 & 0.960 & -- \\
gte-large & 0.791 & $-1.280$ & 0.919 & 0.961 & -- \\
bge-small-en-v1.5 & 0.791 & $-1.882$ & 0.776 & 0.918 & -- \\
multilingual-e5-large & 0.798 & $-2.299$ & 0.873 & 0.958 & -- \\
bge-base-en-v1.5 & 0.799 & $-1.860$ & 0.762 & 0.909 & -- \\
all-distilroberta-v1 & 0.801 & $-1.856$ & 0.629 & 0.858 & -- \\
all-mpnet-base-v2 & 0.805 & $-2.456$ & 0.545 & 0.860 & -- \\
bge-large-en-v1.5 & 0.812 & $-2.108$ & 0.751 & 0.914 & -- \\
\bottomrule
\end{tabular}
\end{center}

\subsection{FEVER: 9 Models}

\begin{center}
\small
\begin{tabular}{@{}lcccc@{}}
\toprule
\textbf{Model} & \textbf{MTEB $\rho$} & \textbf{Sep.} & \textbf{AUROC} & $\bar{s}_S$ / $\bar{s}_R$ \\
\midrule
all-MiniLM-L6-v2 & 0.772 & 0.786 & 0.712 & 0.686 / 0.590 \\
e5-base-v2 & 0.775 & 0.758 & 0.708 & 0.868 / 0.840 \\
nomic-embed-v1.5 & 0.780 & 0.689 & 0.689 & 0.811 / 0.762 \\
gte-base & 0.783 & 0.680 & 0.688 & 0.920 / 0.900 \\
e5-small-v2 & 0.787 & 0.852 & 0.727 & 0.881 / 0.850 \\
e5-large-v2 & 0.787 & 0.731 & 0.697 & 0.855 / 0.827 \\
bge-small-en-v1.5 & 0.791 & 0.917 & 0.748 & 0.799 / 0.730 \\
bge-base-en-v1.5 & 0.799 & 0.949 & 0.750 & 0.775 / 0.701 \\
all-mpnet-base-v2 & 0.805 & 0.815 & 0.723 & 0.688 / 0.582 \\
\bottomrule
\end{tabular}
\end{center}

\subsection{SummEval: 7 Models}

\begin{center}
\small
\begin{tabular}{@{}lccc@{}}
\toprule
\textbf{Model} & \textbf{MTEB $\rho$} & \textbf{Mean Spearman} & \textbf{Pairwise Acc.} \\
\midrule
GloVe 6B 300d & 0.554 & 0.196 & 0.573 \\
paraphrase-MiniLM-L3 & 0.762 & 0.287 & 0.610 \\
all-MiniLM-L6-v2 & 0.772 & 0.247 & 0.593 \\
e5-small-v2 & 0.787 & 0.272 & 0.602 \\
bge-small-en-v1.5 & 0.791 & 0.308 & 0.618 \\
bge-base-en-v1.5 & 0.799 & 0.330 & 0.625 \\
all-distilroberta-v1 & 0.801 & 0.254 & 0.597 \\
\bottomrule
\end{tabular}
\end{center}

\end{document}
